%%%%%%%%%%%%%%%%%%%%%%%%%%%%%%%%%%%%%%%%%%%%%%%%%%%%%%%%%%%%%%%%%%%%%%%%%%%%%
% APPENDIX: DETAILED THEORETICAL MODEL
%%%%%%%%%%%%%%%%%%%%%%%%%%%%%%%%%%%%%%%%%%%%%%%%%%%%%%%%%%%%%%%%%%%%%%%%%%%%%

\appendix
\section{Detailed Theoretical Model}
This appendix provides additional details and derivations related to the theoretical model presented in the main text.

\subsection{Supplier Entry and Auction Environment}
Assume that each supplier \(i\) has a constant marginal cost \(c\) and enters the auction with probability \(p_O\) in ordinary procurement and \(p_U\) in urgent procurement, where \(p_U < p_O\). With \(N\) potential suppliers, the expected number of entrants is:
\[
E[N_O] = N \cdot p_O, \quad E[N_U] = N \cdot p_U.
\]
The reduction in supplier entry under urgency is a key factor driving the higher equilibrium prices.

\subsection{Derivation of Equilibrium Prices}
Consider the bidding behavior under the three procurement regimes:
\begin{itemize}
    \item \textbf{Ordinary Procurement:} With numerous bidders, competitive bidding drives the equilibrium bid close to the marginal cost:
    \[
    P_O = \min \{b_i^{O}\} \approx c.
    \]
    \item \textbf{Urgent Procurement:} Reduced competition results in a markup \(\delta\):
    \[
    P_U = c + \delta.
    \]
    \item \textbf{Litigated Procurement:} Judicial enforcement introduces an extra cost:
    \[
    P_{Litigated} = c + \delta + \pi \cdot \gamma.
    \]
\end{itemize}

\subsection{Welfare Implications}
Let \(V\) denote the public valuation per unit. Social welfare under each regime is:
\[
EW_j = Q_j \cdot (V - P_j), \quad j \in \{O, U, \text{Litigated}\}.
\]
The comparative statics show that increased urgency (higher \(\delta\)) and stronger judicial penalties (higher \(\pi\) or \(\gamma\)) reduce welfare.

\subsection{Comparative Statics}
The model permits the analysis of changes in key parameters:
\begin{itemize}
    \item An increase in \(\delta\) (reflecting higher urgency) increases \(P_U\) and reduces \(EW_U\).
    \item Higher judicial intervention (\(\pi\)) or penalty severity (\(\gamma\)) further elevates \(P_{Litigated}\) and intensifies welfare losses.
\end{itemize}

\subsection{Proofs of Theorems}
For completeness, we outline the key steps of the proofs:

\begin{proof}[Proof of Theorem 1]
    \begin{enumerate}
        \item In ordinary procurement, competitive entry drives \(P_O \approx c\).
        \item In urgent procurement, the markup \(\delta\) ensures \(P_U = c + \delta > P_O\).
        \item In litigated procurement, the additional term \(\pi \cdot \gamma\) leads to \(P_{Litigated} = c + \delta + \pi \cdot \gamma > P_U\).
    \end{enumerate}
    Thus, \(\mathbb{E}[P_{Litigated}] > \mathbb{E}[P_U] > \mathbb{E}[P_O]\).
\end{proof}

\begin{proof}[Proof of Theorem 2]
    The welfare loss under litigated procurement is:
    \[
    \Delta W_{Litigated} = Q_U \cdot \left((P_{Litigated} - P_O) + \pi \cdot \gamma\right),
    \]
    which follows directly from the definitions of \(EW_O\) and \(EW_{Litigated}\). \(\Box\)
\end{proof}
